\documentclass[a4paper,14pt]{extarticle}
\usepackage{array,amsmath,fancybox,amsfonts,graphicx,amsthm,amssymb,wasysym,latexsym,slashed,anysize,subfigure,calc}
\usepackage[spanish]{babel}
\usepackage[utf8]{inputenc}
\usepackage{fancyhdr,float,hyperref}
\usepackage{lettrine}


%lewis y química
\usepackage{chemformula,elements}
%chemformula, provides \ch{} with the !(<below>)(<formula>) syntax and \chlewis;
%elements, provides \elconf and \writeelconf.
%ex:
%\ch{
%  !(\elconf{Li})( "\chlewis{180.}{Li}" ) +
%  !(\elconf{F})( "\chlewis{0.90:180:270:}{F}" )
%   ->
%  !(\writeelconf{2})( Li+ ) +
%  !(\writeelconf{2,2+6})( "\chlewis{0:90:180:270:}{F}" {}- )
%}


%entorno para química: texto y math mode
%\newcommand*\chem[1]{\textup{#1}}
\newcommand*\chem[1]{\ensuremath{\mathrm{#1}}}

%column
\usepackage{multicol}
%uso: \begin{multicols}{3}

%enumitem: configurar listas
\usepackage{enumitem}

%para los entornos de soluciones
\usepackage{mdframed,xcolor,color}
	%uso: 	entorno (begin,end): {mdframed}[backgroundcolor=brown!10] 

%subrayado y highlight. Tiene que estar detrás de usepackage[color]
\usepackage{soul}
	%\textst{normal} 
	\DeclareRobustCommand{\hlcyan}[1]{{\sethlcolor{cyan}\hl{#1}}} %en cyan!
		%\hl{highlighted text} %uso


%tamaños fuentes extra: 8,9,14,17,20
	\usepackage{extsizes}
	
 	%fuente librebaskerville
 	\usepackage[T1]{fontenc}
	\usepackage{librebaskerville}
	
	%fracciones $$ en texto más elaboradas
	 \usepackage{xfrac}
	 
	%permitir que una ecuación larga esté separada por un salto de página
	\allowdisplaybreaks[4]
		
	%gráficos varios 
	\usepackage{graphics}
	
	%fancy tables
	\usepackage{booktabs}
	\usepackage[small,centerlast,up]{caption}

	
	%mejora en la exposición de guarismos
	\usepackage{numprint}
	
\usepackage{everypage,draftwatermark} %Marca de agua
\SetWatermarkText{\vspace*{2cm}

\hspace*{2cm}\parbox{2\textwidth}{\centering DIST\\B1-FyQ}}
\SetWatermarkScale{1.5}
\SetWatermarkLightness{0.9}

\DeclareMathOperator{\lin}{lin}
\DeclareMathOperator{\id}{\textbf{id}}
\DeclareMathOperator{\modulo}{mod}

\newcommand{\mc}[3]{\multicolumn{#1}{#2}{#3}}
\newcommand{\halmos}{
\vspace*{-3ex}

\begin{flushright}
\rule{1mm}{2.5mm} \end{flushright}
}
\renewcommand{\qed}{\halmos}


%versores: %uso: \uveci\ne\hat{\imath}_{\uvec{i}}+\uvec{j}+\uvec{k}$
\usepackage{bm}
\newcommand{\uveci}{{\bm{\hat{\textnormal{\bfseries\i}}}}}
\newcommand{\uvecj}{{\bm{\hat{\textnormal{\bfseries\j}}}}}
\DeclareRobustCommand{\uvec}[1]{{%
  \ifcsname uvec#1\endcsname
     \csname uvec#1\endcsname
   \else
    \bm{\hat{\mathbf{#1}}}%
   \fi
}}


%algún nuevo operador
\DeclareMathOperator{\cotan}{cotan}


%para textbf en entorno enumerate
\renewcommand{\labelenumi}{%
 \textbf{\theenumi}.
}

%a4
\marginsize{1cm}{1cm}{1cm}{1cm}%requiere el paquete anysize,izq,dcha,arriba,abajo

%a5
%\marginsize{1cm}{1cm}{1cm}{.5cm}%requiere el paquete anysize,izq,dcha,arriba,abajo

\newtheorem{teo}{{\sc \textbf{Teorema}}}
\newtheorem{defi}{{\sc \textbf{Definición}}}
\newtheorem{coro}{{\sc \textbf{Corolario}}}


%opening
\title{\textbf{	}}
\date{}
\pagestyle{fancy}
\fancyhead[LO,RE]{{\small DOSSIER DE PROBLEMAS}}
\fancyhead[RO,LE]{{\small B1-Dist FyQ}}
%\fancyhead[CO,CE]{Ricardo Torres\\http://tinyurl.com/iesCMGfyq}



\begin{document}
\pagenumbering{gobble}

%	\section*{Ejercicios propuestos Tema 3: óptica ondulatoria}

%\begin{center}
%	\hl{2, 4, 10 p\'ag. 241}\hspace{2cm}\hl{24 p\'ag. 242}
%	\end{center}

%\vspace*{.1cm}
%
%	\section*{Ejercicios adicionales sobre la luz}
%


\begin{enumerate}
  \setlength\itemsep{1.5em}
    
    
    \item \textbf{(2 PUNTOS)} 
    
    Formula y nombra:
    	%\begin{multicols}{2}
    		\begin{enumerate}[label=\alph*)]
    			\item Los ácidos ternarios de los tres primeros elementos del grupo 15.
    			\begin{mdframed}[backgroundcolor=brown!10]  
    			Se trata de nombrar los ácidos que forman el nitrógeno, el fósforo y el arsénico. Evidentemente habrá que tener cuidado con el fósforo y el arsénico, que puede formar dos tipos de ácidos ternarios.
    			\end{mdframed}
    			\item Las sales que forma el azufre con el segundo y tercer elemento del grupo cuyo último electrón corresponda a un estado $ns^2$.
    			\begin{mdframed}[backgroundcolor=brown!10]  
    			El grupo cuyo último electrón es un $ns^2$ es evidentemente el grupo $2$. De ahí que haya que seleccionar las sales que forma el azufre con magnesio y calcio: $\chem{MgS}$ y $\chem{CaS}$.
    			\end{mdframed}
    			\item Los peróxidos que se forman con los metales alcalinos correspondientes a los periodos $3$ y $4$.
    			\begin{mdframed}[backgroundcolor=brown!10]  
    			Los metales alcalinos corresponden al grupo $1$, de modo que los elementos que buscamos son sodio y potasio. Formular sus peróxidos es poco más que trivial.
				\end{mdframed}    			 
    			\item Las sales ternarias que forma el hierro con los aniones resultantes de la disociación completa de los ácidos ternarios correspondientes a los elementos con configuración $\elconf{se}$ y $\elconf{i}$.
    			\begin{mdframed}[backgroundcolor=brown!10]  
    			Los elementos considerados son el selenio y el yodo. Sus ácidos ternarios son fáciles de encontrar y permiten, por disociación (emisión de iones $\chem{H^+}$) formar las sales que se buscan.
    			\end{mdframed}
    		\end{enumerate}
    		\qed
    	%\end{multicols}
    	\newpage
    	
    \item \textbf{(1 PUNTO)}
    
    	\begin{enumerate}[label=\alph*)]
    		\item ¿Cuál es la diferencia entre un alcano y un alquino?
    		\begin{mdframed}[backgroundcolor=brown!10]  
    		Un alcano es un hidrocarburo lineal de cadena simple, mientras que un alquino incluye uno o más enlaces triples.
    		\end{mdframed}
    		\item  Nombra y formula (en forma semidesarrollada) los seis primeros alcanos. \textit{Idem} para los cuatro primeros alquenos en todas las formas \textit{distintas} en que sea posible que contengan dobles enlaces.
    		\begin{mdframed}[backgroundcolor=brown!10]  
    		Los alcanos pueden consultarse en cualquier manual. Los alquenos son algo más complejos porque es necesario pensar las posibles maneras de colocar dobles enlaces.
    		\end{mdframed}
    	\end{enumerate} 
    \qed
    
    
    \item \textbf{(1 PUNTO)}
    
     A la vista del proceso que hemos analizado para la formación de sales ternarias a partir la disociación de ácidos oxoácidos, ¿qué relación existe entre un ácido carboxílico y un éster? 
     \begin{mdframed}[backgroundcolor=brown!10]  
     La pregunta casi incluye la respuesta: evidentemente tendrá que ver con la disociación de ácidos ternarios. En efecto, echando un vistazo a lo que hemos visto en clase respecto a la forma del grupo carboxilo y el grupo éster, observamos que los ésteres pueden entenderse \textit{como las sales a las que da lugar un ácido orgánico disociado.}
     \end{mdframed}
    \qed
    
    \newpage
    
    \item \textbf{(3 PUNTOS)} 
    
    Considera los elementos desconocidos $^{8}\alpha$, $^{18}\beta$, $^{34}\delta$ y $^{38}\epsilon.$ 
    	\begin{enumerate}[label=\alph*)]
    		\item Escribe las configuraciones electrónicas de los cuatro elementos. 
    		\begin{mdframed}[backgroundcolor=brown!10]  
    				\begin{itemize}
    					\item $[^8\alpha]=\elconf{O}$
    					\item $[^{18}\beta]=\elconf{Ar}$
    					\item $[^{34}\delta]=\elconf{Se}$
    					\item $[^{38}\epsilon]=\elconf{Sr}$
    				\end{itemize}
    			\end{mdframed}
    		\item Justifica el tipo de enlace y la fórmula del compuesto que forman los elementos $\delta$ y $\epsilon$.
    		\begin{mdframed}[backgroundcolor=brown!10]  
    		A la vista de la configuración $\delta$ es un no metal del grupo $16$ que \textit{busca} ganar dos electrones para cerrar su capa de valencia.   		
    		 Por otro lado $\epsilon$ es un metal alcalinotérreo, electropositivo, y que \textit{busca} perder dos electrones para encontrar la estabilidad. 	Se unirán por enlace iónico dando lugar al compuesto $\epsilon\delta$.
    		 
    		 % Como compuesto iónico probablemente será un sólido a temperatura ambiente debido a la fuerte interacción de carácter electromagnético que actúa interatómicamente en la red cristalina. Será soluble en disolventes polares y conducirá la electricidad disuelto o fundido, estados en que dispondrá de cargas móviles.
    		\end{mdframed}
    		\item Justifica la existencia de la molécula $\alpha_2$ e indica el tipo de enlace que se produce.
    		\begin{mdframed}[backgroundcolor=brown!10]  
    		$\alpha$ es un no metal del grupo $16$. Para conseguir la estabilidad puede interaccionar con un átomo de la misma especie compartiendo un par de electrones (par de \textit{Lewis}) de modo que consiga alcanzar una estructura electrónica de gas noble.
    		\end{mdframed}
    		\item Indica el tipo de enlace que formarán los elementos $\beta$ y $\alpha$.
    		\begin{mdframed}[backgroundcolor=brown!10]  
    		$\beta$ es un gas noble y no forma compuestos.
    		\end{mdframed}
    	\end{enumerate}
   
    \qed
    	
    
    \newpage
    
    \item \textbf{(3 PUNTOS)} 
    
    El modelo de Bohr para el átomo de hidrógeno permite explicar la existencia de un espectro atómico discreto debido a un mecanismo de transición entre órbitas. En efecto, el único electrón del hidrógeno es capaz de pasar de una órbita más externa a una más interna emitiendo luz; o pasar de una órbita más interna a una más externa absorbiendo luz.
    
    Según el modelo de Bohr, la energía de un electrón en una órbita descrita por el número cuántico $n$ puede escribirse como
    	$$E_n=-\frac{13.6}{n^2},$$
    medida en $eV$, una unidad de energía que equivale aproximadamente a $1.6\cdot 10^{-19}\;J$.
    		
    	\begin{enumerate}[label=\alph*)]
    		\item La cantidad de energía luminosa emitida o absorbida en la transición de un electrón desde un nivel $n_1$ a un nivel $n_2$ es $$\Delta E=|E_{n_1}-E_{n_2}|.$$
    		
    	 Calcula las cuatro energías, en $eV$, que puede emitir un electrón cuando \textit{cae} al nivel fundamental $n=1$ procedente de las órbitas correspondientes a $n=2$, $n=3$, $n=4$ y $n=5$.
    	 \begin{mdframed}[backgroundcolor=brown!10]  
    	 Las cuatro transiciones llevan aparejadas una energía dada por la diferencia, en valor absoluto, entre las energías de cada nivel según el modelo de Bohr. El cálculo, en $eV$, es trivial.
    	 \end{mdframed}
    		\item Planck propuso que la energía luminosa emitida o absorbida en una transición correspondía a la emisión o absorción de una partícula elemental llamada fotón que podía caracterizarse por una frecuencia $f$, medida en $s^{-1}$, según la expresión 
    		$$\Delta E=hf,$$
    		donde $h\simeq 6.63\cdot 10^{-34}\;J\cdot s$ es la llamada constante de Planck. 
    		Calcula la frecuencia asociada a los fotones correspondientes a las transiciones del apartado a).
    		\begin{mdframed}[backgroundcolor=brown!10]  
    		Trivial. Lo único que hay que tener en cuenta es que para usar la ecuación de Planck hay que usar la constante homónima, de modo que tendremos que ser cuidadosos con las unidades si queremos relacionar las energías asociadas a las transiciones del ejercicio a), originalmente en $eV$. 
    		\end{mdframed}
    		\item Años antes Rydberg, de forma completamente experimental, había llegado a la conclusión de que la frecuencia de un fotón emitido como resultado de la transición de un electrón desde el nivel $n_i$ a un nivel inferior $n_f$ podía escribirse como 
    		$$f=\beta\left(\frac{1}{n_f^2}-\frac{1}{n_i^2}\right),$$
    		donde $\beta$ es una cierta constante. Calcula el valor de $\beta$ y sus unidades a partir de la información de los apartados a) y b).
    		\begin{mdframed}[backgroundcolor=brown!10]  
    		Tan solo hay que ser cuidadoso con las unidades (¡toda cantidad física debe llevarlas so pena de estar hablando sobre nada!). Elige una transición de las que has calculado en los apartados a) y b), y calcula $\beta$ introduciendo los valores de la frecuencia y los números cuánticos de partida y de llegada correspondientes. 
    		
    		Evidentemente $\beta$ debe ser positivo, puesto que $n_i>n_f$, y debe tener las mismas unidades que $f$ debido a que los números cuánticos son adimensionales.
    		\end{mdframed}
    	\end{enumerate}
		\qed    	

\end{enumerate}

\end{document}

